\chapter{Исследование предметной области} \label{ch1}

Данная глава посвящена исследованию предметной области. 

В параграфе \ref{ch1:sec1} рассматривается актуальность применения СУБД Greenplum в современных аналитических системах, а также обосновывается необходимость разработки системы мониторинга. Параграф \ref{ch1:sec2} посвящён анализу существующих решений для мониторинга Greenplum, в частности коммерческой системы Greenplum Command Center, и выявлению её основных функциональных ограничений. В параграфе \ref{ch1:sec3} рассматриваются теоретические основы мониторинга СУБД, а также особенности мониторинга Greenplum, включая рекомендации разработчиков по проведению основных мониторинговых процедур. Параграф \ref{ch1:sec4} содержит формирование требований к разрабатываемой системе мониторинга на основе проведённого анализа предметной области и существующих решений.

\section{Актуальность проблемы} \label{ch1:sec1}

В настоящее время СУБД Greenplum получила достаточно широкое распространение в качестве основного хранилища данных и аналитических систем. Это обусловлено несколькими факторами.

Greenplum построен на архитектуре массово-параллельной обработки (MPP), изображенной на \firef{fig:mpp_architecture}. Это означает, что выполнение запросов распределяется между несколькими вычислительными узлами, что позволяет значительно увеличить скорость обработки данных и обеспечить горизонтальную масштабируемость. 

\begin{figure}[ht!] 
	\center
	\includegraphics [scale=0.3] {/Users/vadimvasilev/Study/8 семестр/НИР/latex/SPbPU-student-thesis-template/my_folder/images//1}
	\caption{MPP Архитектура} 
	\label{fig:mpp_architecture}  
\end{figure}

За счет распределения данных по сегментам и параллельного выполнения запросов на сегментах Greenplum показывает лучшую производительность на больших объемах данных по сравнению с традиционнами реляционными СУБД, такими как PostgreSQL. 

Использование зеркальных (mirror) сегментов обеспечивает повышение отказоустойчивости системы, поскольку при выходе из строя основного сегмента выполнение операций автоматически переключается на соответствующий резервный сегмент без прерывания работы всего кластера.

Несмотря на высокую производительность, масштабируемость и встроенные механизмы обеспечения отказоустойчивости, эксплуатация Greenplum требует постоянного контроля состояния системы. В условиях промышленной эксплуатации даже кратковременные отклонения в работе СУБД могут привести к снижению общей производительности кластера или возникновению критических ошибок. По этой причине задача мониторинга состояния СУБД является неотъемлемой частью обязанностей администратора базы данных.

На данный момент представлено только одно специализированное решение для мониторинга СУБД Greenplum. Это коммерческое решение -- Greenplum Command Center. Данный продукт является компонентом коммерческого дистрибутива Tanzu VMware Greenplum. Конкретная стоимость дистрибутива в открытом доступе не сообщается. Подробное рассмотрение данной системы будет приведено в слудующем параграфе  \ref{ch1:sec2}.

Доступных на территории Российской Федерации бесплатных комплексных решений, обеспечивающих мониторинг Greenplum в настоящее время не представлено. 

Приведенные факторы подтверждают актуальность и практическую значимость данной работы.

\section{Обзор существующих систем мониторинга Greenplum} \label{ch1:sec2}

Greenplum Command Center является инструментом мониторинга и управления кластером СУБД Greenplum, входящим в состав коммерческого дистрибутива Greenplum. Командный центр позволяет отслеживать показатели производительности системы, анализировать работоспособность кластера, а так же позволяет  выполнять задачи администрирования кластера.

Система предоставляет веб-интерфейс, позволяющий в режиме реального времени отслеживать ключевые системные показатели работы кластера, включая загрузку вычислительных узлов, использование оперативной памяти, нагрузку на процессор, а также активность выполнения SQL-запросов.

В открытых источниках есть примеры интерфейса системы, который содержит несколько дашбордов. Дашборд изображенный на \firef{fig:tanzu_1}, отображает основные показатели производительности кластера. На отдельном дашборде, изображенном на \firef{fig:tanzu_2} доступен мониторинг активности запросов к базе данных.

\begin{figure}[ht!] 
	\center
	\includegraphics [scale=0.5] {/Users/vadimvasilev/Study/8 семестр/ВКР (Диплом)/latex/SPbPU-student-thesis-template/my_folder/images//tanzu_1}
	\caption{Дашборд системных метрик в Greenplum Command Center} 
	\label{fig:tanzu_1}  
\end{figure}

\begin{figure}[ht!] 
	\center
	\includegraphics [scale=0.5] {/Users/vadimvasilev/Study/8 семестр/ВКР (Диплом)/latex/SPbPU-student-thesis-template/my_folder/images//tanzu_2}
	\caption{Дашборд активности запросов в Greenplum Command Center} 
	\label{fig:tanzu_2}  
\end{figure} 

В системе настроен определенный список событий, при возникновении которых пользоватлю отправляется уведомление, посмотреть которое можно в специальной вкладке <<Уведомления>>, пример которой можно увидеть на \firef{fig:tanzu_3}. Получение уведомлений на электронную почту требует отдельной настройки со стороны администратора системы. 

\begin{figure}[ht!] 
	\center
	\includegraphics [scale=0.6] {/Users/vadimvasilev/Study/8 семестр/ВКР (Диплом)/latex/SPbPU-student-thesis-template/my_folder/images//tanzu_3}
	\caption{Пример уведомлений в Greenplum Command Center} 
	\label{fig:tanzu_3}  
\end{figure} 

Несмотря на достаточно широкий функционал, данный инструмент имеет ряд недостатков. В первую очередь он распространяется как часть коммерческого продукта и требует использования дорогостоящей лицензии. Кроме того встроенная система уведомлений без дополнительной настройки позволяет узнавать о критических инцидентах только непосредственно из веб-интерфейса командного центра, что препятствует оперативному реагированию на внештатные ситуации.

Таким образом, Greenplum Command Center представляет собой коммерческое решение для мониторинга СУБД Greenplum, недоступное широкому кругу пользователей, что подтверждает необходимость разработки альтернативных систем, ориентированных на доступность, и простоту эксплуатации. 
 
\section{Особенности мониторинга Greenplum} \label{ch1:sec3} 

Мониторинг СУБД представляет собой процесс сбора и анализа информации о состоянии системы с целью обеспечения её стабильной и эффективной работы. Он включает в себя отслеживание различных показателей, характеризующих использование ресурсов и выполнение операций.

К основным задачам мониторинга относятся \cite{monitoringbook}:

\begin{itemize}
	\item контроль системных показателей, таких как загрузка процессора, использование памяти и дисковых ресурсов;
	\item отслеживание доступности системы;
	\item анализ выполняемых запросов и их длительности;
	\item выявление узких мест в производительности;
	\item контроль ключевых показателей работы СУБД.
\end{itemize}

Реализация мониторинга требует опоры на техническую документацию и регламенты администрирования системы. В связи с этим теоретической основой формирования требований к системе мониторинга в рамках данной работы послужили рекомендации разработчиков Greenplum. В документации изложен перечень регламентных мониторинговых процедур, направленных на обеспечение высокой доступности и стаблильной производительности кластера. 

Основные рекомендованные процедуры мониторинга можно разделить на 4 основных группы:

\begin{enumerate}[1.]
	\item Мониторинг состояния кластера Greenplum
	{\small 
	\begin{longtable}{|p{3.5cm}|p{6.5cm}|p{2.2cm}|p{2.8cm}|}
		\caption{Рекомендованные процедуры проверки состояния кластера Greenplum}
		\label{tab:table_2} \\
		
		\hline
		\textbf{Процедура} & \textbf{Необходимое действие} & \textbf{Степень значимости} & \textbf{Рекомендуемая частота проверки} \\
		\hline
		\endfirsthead
		
		\hline
		\textbf{Процедура} & \textbf{Необходимое действие} & \textbf{Степень значимости} & \textbf{Рекомендуемая частота проверки} \\
		\hline
		\endhead
		
		Проверка работоспособности и доступности мастер-сегмента &
		Необходимо выполнить запрос: \newline
		\texttt{SELECT count(*) FROM gp\_segment\_configuration;} &
		Критическая & Каждые 5-10 минут \\
		\hline
		
		Проверка списка недоступных сегментов &
		Необходимо выполнить запрос: \newline
		\texttt{SELECT * FROM gp\_segment\_configuration WHERE status = 'd';} &
		Высокая & Каждые 5-10 минут \\
		\hline
		
		Проверка активных, но не синхронизированных сегментов &
		Необходимо выполнить запрос: \newline
		\texttt{SELECT * FROM gp\_segment\_configuration WHERE mode = 'n' and status = 'u' and content <> -1;} &
		Высокая & Каждые 5-10 минут \\
		\hline
		
		Проверка распределённого выполнения запроса &
		Необходимо выполнить запрос: \newline
		\texttt{SELECT gp\_segment\_id, count(*) FROM gp\_dist\_random('pg\_class') GROUP BY 1;} \newline
		Для каждого сегмента должна быть возвращена 1 строка &
		Критическая & Каждые 5-10 минут \\
		\hline
		
		Проверка состояния репликации master &
		Необходимо psql команду: \newline
		\texttt{psql <dbname> -c 'SELECT pid, state FROM pg\_stat\_replication;'} \newline
		Значение \texttt{state} должно быть <<STREAMING>> &
		Критическая & Каждые 5-10 минут \\
		\hline
		
	\end{longtable}
	}
	
	\item Мониторинг аппаратного обеспечения и операциинной системы
	
	{\small 
		\begin{longtable}{|p{3.5cm}|p{6.5cm}|p{2.2cm}|p{2.8cm}|}
			\caption{Рекомендованные процедуры проверки аппаратного обеспечения и операционной системы}
			\label{tab:table_3} \\
			
			\hline
			\textbf{Процедура} & \textbf{Необходимое действие} & \textbf{Степень значимости} & \textbf{Рекомендуемая частота проверки} \\
			\hline
			\endfirsthead
			
			\hline
			\textbf{Процедура} & \textbf{Необходимое действие} & \textbf{Степень значимости} & \textbf{Рекомендуемая частота проверки} \\
			\hline
			\endhead
			
			Проверка использования дискового пространства на томах, используемых для хранения данных базы данных Greenplum &
			Необходимо проверить наличие свободного дискогового пространства. Рекомендуется предпринять меры при достижении 75\% занятого объема диска &
			Критическая & Каждые 10-30 минут \\
			\hline
			
			Проверка состояния дисков RAID массива &
			Совместно с администраторами сервера необходимо проверить состояние всех дисков RAID массива. При наличии ошибок или выходе из строя каких-либо аппаратных компонент незамедлительно заменить их &
			Критическая & При создании кластера, или при наличии признаков неисправности \\
			\hline
			
			Проверка состояния сетевых интерфесов &
			Необходимо проверить сетевые интерфейсы на наличие ошибок и потерянных пакетов передачи данных &
			Высокая & Каждый час  \\
			\hline
		\end{longtable}
	}
	
	\item Мониторинг состояния каталогов
	
	{\small 
		\begin{longtable}{|p{3.5cm}|p{6.5cm}|p{2.2cm}|p{2.8cm}|}
			\caption{Рекомендованные процедуры проверки состояния каталогов}
			\label{tab:table_4} \\
			
			\hline
			\textbf{Процедура} & \textbf{Необходимое действие} & \textbf{Степень значимости} & \textbf{Рекомендуемая частота проверки} \\
			\hline
			\endfirsthead
			
			\hline
			\textbf{Процедура} & \textbf{Необходимое действие} & \textbf{Степень значимости} & \textbf{Рекомендуемая частота проверки} \\
			\hline
			\endhead
			
			Проверка согласованности каталогов каждой базы данных &
			Необходимо убедиться, что каталог на каждом узле кластера согласован и находится в исправном состоянии. Для этого рекомендуется использовать утилиту \texttt{gpcheckcat}: \texttt{gpcheckcat -O}  &
			Высокая & Каждую неделю \\
			\hline
			
			Проверка целостности системного каталога &
			Рекомендуется использовать утилиту \texttt{gpcheckcat}: \texttt{gpcheckcat -R pgclass} для выявления несогласованных записей в системных таблицах \texttt{pg\_class} и \texttt{pg\_attribute} &
			Критическая & Каждый месяц \\
			\hline
			
			Проверка корректности схем и временных пространств имен &
			Необходимо убедиться в отсутствии некорректных временных схем и отсутствующих определений пространств имен в системном каталоге. Для этого рекомендуется использовать утилиту \texttt{gpcheckcat}: \texttt{gpcheckcat -R namespace} &
			Высокая & Каждый месяц  \\
			\hline
			
			Проверка зависимостей системного каталога &
			Необходимо убедиться в отсутствии зависимостей от несуществующих объектов в системном каталоге базы данных. Для этого рекомендуется использовать утилиту \texttt{gpcheckcat}: \texttt{gpcheckcat -R dependency} &
			Высокая & Каждый месяц  \\
			\hline
		\end{longtable}
	
	}
	
	\item Мониторинг состояния базы данных
	
	{\small 
		\begin{longtable}{|p{3.5cm}|p{6.5cm}|p{2.2cm}|p{2.8cm}|}
			\caption{Рекомендованные процедуры проверки состояния базы данных}
			\label{tab:table_5} \\
			
			\hline
			\textbf{Процедура} & \textbf{Необходимое действие} & \textbf{Степень значимости} & \textbf{Рекомендуемая частота проверки} \\
			\hline
			\endfirsthead
			
			\hline
			\textbf{Процедура} & \textbf{Необходимое действие} & \textbf{Степень значимости} & \textbf{Рекомендуемая частота проверки} \\
			\hline
			\endhead
			
			Проверка таблиц с отсутствующей статистикой &
			Необходимо проверять таблицы, для которых отсутствует актуальная статистика. Для этого рекомендуется выполнить запрос к представлению \texttt{gp\_toolkit.gp\_stats\_missing}: \texttt{SELECT * FROM gp\_toolkit.gp\_stats\_missing;} &
			Средняя & После загрузки большого объема данных \\
			\hline
			
			Проверка таблиц с избыточным неиспользуемым пространством &
			Необходимо выявлять таблицы, содержащие значительный объем неиспользуемого пространства, не освобождаемого обычной очисткой. Для этого рекомендуется выполнить запрос к представлению \texttt{gp\_toolkit.gp\_bloat\_diag}: \texttt{SELECT * FROM gp\_toolkit.gp\_bloat\_diag;} &
			Средняя & Каждую неделю или каждый месяц \\
			\hline
			
			Очистка таблиц от удаленных строк &
			Необходимо освобождать пространство, занимаемое удаленными строками, для его повторного использования. Для этого рекомендуется выполнять команду \texttt{VACUUM <table>}. &
			Высокая & Каждый день \\
			\hline
			
			Обновление статистики таблиц &
			Необходимо поддерживать статистику таблиц в актуальном состоянии для корректной работы оптимизатора запросов. Для этого рекомендуется использовать утилиту \texttt{analyzedb}: \texttt{analyzedb -d <database> -a}. &
			Высокая & После загрузки большого объема данных \\
			\hline
			
			Резервное копирование данных базы данных &
			Необходимо создавать актуальные резервные копии мастер-узла и сегментов кластера для возможности восстановления базы данных в случае отказа. Для этого рекомендуется использовать утилиту \texttt{gpbackup}. &
			Критическая & Каждый день \\
			\hline
			
			Обслуживание системных каталогов &
			Необходимо выполнять очистку, переиндексацию и обновление статистики системных таблиц для поддержания эффективной работы системного каталога. Для этого рекомендуется использовать команды \texttt{VACUUM}, \texttt{REINDEX SYSTEM} и утилиту \texttt{analyzedb -s pg\_catalog}. &
			Высокая & Каждую неделю \\
			\hline
		\end{longtable}
		
	}
	
\end{enumerate} 

Представленные рекомендации демонстрируют, что обеспечение стабильной работы Greenplum требует контроля большого числа параметров. Однако значительная часть перечисленных процедур относится к периодическому регламентному обслуживанию, предполагает участие администратора и, соответственно, выходит за рамки разрабатываемой системы мониторинга.

Для разработки автоматизированной системы мониторинга наибольший интерес представляют процедуры непрерывного оперативного контроля, позволяющие своевременно фиксировать отклонения в работе СУБД.

В рамках данной работы был выбран подход, основанный на анализе показателей, доступных на мастер-узле кластера. Такой подход позволяет получить обобщенное представление о текущем состоянии системы, включая сведения о нагрузке, доступности сегментов и возникновении критических ошибок, а также сокращает сложность архитектуры разрабатываемого решения, поскольку отпадает необходимость установки дополнительных компонентов на каждый вычислительный сегмент. При этом доступных на управляющем узле данных достаточно для выявления большинства критических ситуаций, связанных с производительностью, активностью и доступностью СУБД. Таким образом, выбранный подход обеспечивает баланс между полнотой собираемой диагностической информации и простотой практической реализации системы мониторинга.

\section{Формирование требований к разрабатываемой системе} \label{ch1:sec4}

На основании проведённого анализа существующего решения Greenplum Command Center, а также рассмотренных особенностей мониторинга СУБД Greenplum, можно сформулировать основные требования к разрабатываемой системе мониторинга:

\begin{itemize}
	\item обеспечение централизованного сбора и хранения диагностической информации о состоянии кластера. К числу контролируемых параметров должны относиться системные метрики вычислительного узла, показатели доступности сегментов, сведения о нагрузке на базу данных, а также информация о возникновении критических ошибок в работе СУБД;
	\item наличие информативного дашборда, отображающего основные показатели СУБД;
	\item наличие механизма автоматического оповещения посредствам электронной почты о критических инцидентах в работе СУБД;
	\item обеспечение доступности и простоты внедрения разрабатываемого решения. Система не должна требовать использования коммерческих лицензий, установки большого числа вспомогательных компонентов или сложной процедуры интеграции в существующую инфраструктуру.
\end{itemize}

\section{Выводы по главе} \label{ch1:conclusion}

В рамках данной главы была рассмотрена предметная область мониторинга СУБД Greenplum. Проведен анализ существующего решения Greenplum Command Center, рассмотрены его основные функциональные возможности и выявлены ключевые недостатки.

Также были рассмотрены теоретические аспекты мониторинга СУБД. На основе которых сформулированы требования к разрабатываемой системе мониторинга, определяющие её функциональные возможности.

Полученные результаты составляют теоретическую и практическую основу для дальнейшей разработки архитектуры и реализации системы мониторинга Greenplum.

\newpage
